

\begin{abstract}

The present work is part of a graduation project carried out within the company \studyDepartment in order to obtain the national diploma of engineer at the \ITBS.

This project aims to design and implement a Business Intelligence and Machine Learning platform dedicated to customer risk assessment and management at the Arab Tunisian Bank (ATB). The platform was built across three successive releases. The first established a secure, role-based environment and a fully automated data orchestration engine that refreshes all banking data every night without human intervention. The second built the analytical foundation, transforming raw operational data into a structured data warehouse and delivering interactive dashboards covering customer portfolio analytics, account monitoring, and a transparent multi-pillar risk scoring system. The third introduced a predictive layer through a supervised machine learning classifier that screens customers for risk, with model tracking, explainability features, and seamless integration into the existing dashboard. The outcome is a unified decision-support platform that gives risk managers, compliance officers, and general management a shared, reliable view of their customer portfolio. \\

\keywordss{Business Intelligence, Machine Learning, Data Warehouse, Customer Risk Assessment, Apache Airflow, Streamlit, MLflow, SHAP, Banking Analytics} \\

\vspace{3cm}
\centerline{\bf Résumé}
\vspace{0.6cm}

Ce travail s'inscrit dans le cadre du projet de fin d'études réalisé au sein de \studyDepartment en vue de l'obtention du diplôme national d'ingénieur à l'\ITBS.

Ce projet a pour objectif de concevoir et de mettre en œuvre une plateforme de Business Intelligence et de Machine Learning dédiée à l'évaluation et à la gestion du risque client au sein de l'Arab Tunisian Bank (ATB). La plateforme a été livrée en trois releases successives. La première a mis en place un environnement sécurisé basé sur les rôles ainsi qu'un moteur d'orchestration entièrement automatisé qui actualise l'ensemble des données bancaires chaque nuit sans intervention humaine. La deuxième a construit la couche analytique, transformant les données opérationnelles brutes en un entrepôt de données structuré et alimentant des tableaux de bord interactifs couvrant l'analyse du portefeuille client, le suivi des comptes et un système de scoring de risque multi-piliers. La troisième a introduit une couche prédictive reposant sur un classificateur supervisé d'apprentissage automatique, accompagné d'un suivi des expériences, d'une explicabilité des prédictions et d'une intégration fluide dans le tableau de bord existant. Le résultat est une plateforme d'aide à la décision unifiée, offrant aux gestionnaires de risque, aux responsables de conformité et à la Direction Générale une vision partagée et fiable de leur portefeuille client. \\

\keywords{Business Intelligence, Machine Learning, Entrepôt de Données, Évaluation du Risque Client, Apache Airflow, Streamlit, MLflow, SHAP, Analytique Bancaire}

\end{abstract}
